% /root/monodromy, 19 November 2026; independently audited by /root.
\subsection{The exact fixed quotient on products of ordinary elliptic curves}
\label{subsec:product-exact-fixed-quotient}
\status{Proved; independent parent proof and source audit passed.
This determines all invariant covectors, not the full group on their kernel.}

\begin{theorem}\label{thm:product-exact-fixed-covectors}
Let $E_1,\ldots,E_n$ be ordinary elliptic curves over a
perfect field of characteristic $p$, where $n\geq2$, and set
\[
 X=\prod_{i=1}^nE_i,\qquad
 L=\boxtimes_{i=1}^n\mathcal O_{E_i}(a_iO_i),\quad a_i\geq3,
 \qquad h_d=d^n\prod_i a_i.
\]
Use the product of the indicated complete embeddings and
the Segre embedding as the fixed projective embedding of $X$.
Then the generic relative Hasse--Witt rank is
\begin{equation}\label{eq:product-elliptic-generic-rank}
 r_d=h_d-1\qquad\text{for every }d\geq1.
\end{equation}
In particular the generic relative coherent Frobenius is
invertible in every degree.

Suppose in addition that $p\mid a_i d$ for every $i$.
For geometric monodromy on its maximal-rank open one has
\begin{equation}\label{eq:product-entire-fixed-space}
 \dim_{\mathbf F_p}(\mathcal V_{d,\bar\eta}^{\vee})^{G_{\rm geom}}
 =\begin{cases}
       2^n-1,&p=2,\\
       n,&p\text{ odd}.
   \end{cases}
\end{equation}
After a finite extension splitting the Frobenius kernel
and normalizing its invariant top Cartier differential,
the arithmetic invariant space is the same space.
Thus these numbers give the dimension of the entire
trivial quotient, rather than a lower bound for it.
\end{theorem}

\begin{proof}
First we prove generic ordinarity.  For each factor
choose a section of $\mathcal O_{E_i}(a_i dO_i)$ with
distinct zeros.  It can be chosen as a product of $d$
sections of $\mathcal O_{E_i}(a_iO_i)$ with disjoint
reduced zero sets.  The resulting product section on
$X$ cuts out an SNC union of coordinate fibres.
Every nonempty stratum is a product of ordinary elliptic
curves, so Frobenius is bijective on every coherent
cohomology group of that stratum.  In the structure-sheaf
Cech spectral sequence, Frobenius is consequently
bijective on every page and on the total cohomology.

For every nonzero defining section, the Kunneth formula
gives $H^{n-1}(X,L^{-d})=0$ and
$h^n(X,L^{-d})=h_d$.  Also $h^n(X,\mathcal O_X)=1$.
The multiplication map to $H^n(X,\mathcal O_X)$ is
surjective because its cokernel lies in
$H^n(Y,\mathcal O_Y)=0$.  Hence the relative coherent
rank is constantly $h_d-1$ on the whole projective
linear system.  Invertibility at the SNC section is
an open condition, and the smooth open is nonempty.
This proves \eqref{eq:product-elliptic-generic-rank}.

The ambient hypersurface family is the same complete
linear system up to a surjective linear projection.
For completeness, a line bundle $N$ of degree at least
three on an elliptic curve is normally generated.
For multiplication $H^0(N)\otimes H^0(N^j)\to H^0(N^{j+1})$
with $j\geq2$, a base-point-free pencil in $H^0(N)$
and $H^1(N^{j-1})=0$ prove surjectivity.
For $j=1$, choose a point $q$ and a base-point-free
pencil in $H^0(N(-q))$.  The pencil exact sequence
after tensoring with $N$ has kernel $\mathcal O(q)$,
whose $H^1$ vanishes, so its products span
$H^0(N^2(-q))$.  The square of a section nonzero at
$q$ spans the remaining one-dimensional quotient.
This proves normal generation.  Tensoring the factors
and grouping Segre monomials proves the required
surjectivity for $X$ in every degree.

We now work over the algebraic closure.  Put
$W=H^0(X,L^d)$ and $P=\mathbf P(W)$, and normalize an
invariant top form $\omega$ by $C(\omega)=\omega$.
The coefficient can be normalized since $X$ is ordinary
and the base field is algebraically closed.
The relative coherent bundle fits in the exact sequence
\begin{equation}\label{eq:product-global-bundle}
 0\longrightarrow\mathcal E_d\longrightarrow
       H^n(X,L^{-d})\otimes\mathcal O_P(-1)
       \longrightarrow\mathcal O_P\longrightarrow0.
\end{equation}
It is a vector bundle on all of $P$, not only on the
ordinary or smooth locus.  Frobenius extends there by
$F_s=s^{p-1}F_X$.

Any geometric monodromy-invariant covector corresponds,
on the invertibility open, to a Frobenius-equivariant
map $\lambda:\mathcal E_d\to\mathcal O$.
It extends over all of $P$.  Indeed, in a regular local
frame its equation is $A^t\lambda=\lambda^{[p]}$ with
$A$ regular.  If its largest pole along any prime
divisor had order $m>0$, the right side would have
a coordinate of valuation $-pm$, whereas every coordinate
on the left has valuation at least $-m$, a contradiction.
There are no divisorial poles; normality of $P$ and local
freeness extend the section also across codimension two.

Dualizing \eqref{eq:product-global-bundle} and taking
global sections, using $H^1(P,\mathcal O_P)=0$, shows
that such a covector is represented by
\begin{equation}\label{eq:product-linear-lift}
 w_s=D(s)\omega,\qquad D\in\End(W),
\end{equation}
modulo scalar multiples of the identity.
Its Frobenius defect lies in the constant one-dimensional
kernel in the Frobenius pullback of the dual sequence.
It is therefore $c s^{[p]}$ for a constant $c$.
Adding $\alpha\operatorname{id}_W$ to $D$ changes this
constant by $\alpha-\alpha^p$; an Artin--Schreier equation
over the algebraically closed field removes it.
Thus every invariant covector lifts to a $D$ satisfying
\begin{equation}\label{eq:product-exact-Cartier-operator}
 C\bigl(s^{p-1}D(s)\omega\bigr)=D(s)\omega
            \quad\text{for all }s\in W.
\end{equation}
Two such exact lifts differ by $c\operatorname{id}$
with $c\in\mathbf F_p$.

The divisibility assumptions give
$L^d=F^*M$ with
$M=\boxtimes_i\mathcal O_{E_i^{(p)}}((a_i d/p)O_i^{(p)})$.
Since $\ker F\simeq\mu_p^n$, its character algebra gives
\begin{equation}\label{eq:product-character-W}
 W=\bigoplus_{\chi\in H}W_\chi,
 \qquad H\simeq\mathbf F_p^n,
 \qquad W_\chi=H^0(X^{(p)},M\otimes\mathcal N_\chi).
\end{equation}
Every $W_\chi$ is nonzero: each factor is a positive
degree line bundle on an elliptic curve.
The normalized Cartier trace is projection to the
trivial character followed by taking the Frobenius
root.  Nontrivial character terms are killed.

Apply \eqref{eq:product-exact-Cartier-operator} to a
section $s\in W_\chi$, and put
$T=\operatorname{pr}_\chi\circ D|_{W_\chi}$.
Only this component contributes to Cartier.  Raising
the equation to the $p$-th power gives
\[
 D(s)^p=s^{p-1}T(s).
\]
We do not assume that the Cartier root preserves the
character space.  Instead choose a nonzero eigenvector
$s_0$ of the finite-dimensional operator $T$, say
$T(s_0)=a s_0$.  The displayed equation implies
$D(s_0)=a^{1/p}s_0$; projecting to $W_\chi$ then gives
$a=a^{1/p}$, so $a\in\mathbf F_p$.
For arbitrary $s\in W_\chi$, substitute $s_0+z s$
and compare the coefficient of $z$.
The left side has no such coefficient, and the right
side gives $s_0^{p-1}(T(s)-a s)=0$.
Since $X$ is integral, $T(s)=a s$.
The same equation now gives $D(s)^p=a s^p$, hence
$D(s)=a s$.  Thus $D$ does preserve every character
space and acts there by a scalar in $\mathbf F_p$:
\begin{equation}\label{eq:product-character-scalar}
 D|_{W_\chi}=c_\chi\operatorname{id},\qquad c_\chi\in\mathbf F_p.
\end{equation}

It remains to determine the allowed functions
$c:H\to\mathbf F_p$.
Choose one nonzero section in each character space
and vary its coefficient independently.
For a nonpure multi-index $(m_\chi)$ with
$\sum m_\chi=p$, the coefficient of the corresponding
monomial in $s^{p-1}D(s)$ is the nonzero factor
$(p-1)!/\prod m_\chi!$ times
$\sum m_\chi c_\chi$.
Cartier kills that monomial unless
$\sum m_\chi\chi=0$ in $H$; if this weight is zero,
the product of the chosen nonzero sections is nonzero
and its Frobenius root is nonzero.
Independence of the coefficient monomials therefore
shows that the exact conditions are
\begin{equation}\label{eq:product-scalar-relations}
 \sum_\chi m_\chi c_\chi=0
 \quad\text{whenever}\quad
 \sum m_\chi=p,\quad\sum m_\chi\chi=0,
 \quad(m_\chi)\text{ is not pure}.
\end{equation}
Pure monomials already give the right side of
\eqref{eq:product-exact-Cartier-operator}.

When $p=2$, every zero-weight pair consists of two
equal characters, so there are no conditions in
\eqref{eq:product-scalar-relations}.  The space of
exact lifts has dimension $2^n$ over $\mathbf F_2$.
When $p$ is odd, put $a_x=c_x-c_0$.
The multiset consisting of $p-2$ zeros, $x$, and $-x$
gives $a_{-x}=-a_x$.
The multiset consisting of $p-3$ zeros, $x$, $y$, and
$-x-y$ then gives $a_{x+y}=a_x+a_y$.
The only nontrivial pure exception is $p=3$, $x=y$;
there the same additivity follows already from
$a_{-x}=-a_x$, since $2x=-x$.
Thus $a$ is an $\mathbf F_p$-linear functional on $H$.
Conversely every affine-linear function
$c_x=c_0+\ell(x)$ satisfies all displayed relations.
The space of exact lifts has dimension $n+1$.

Finally quotient the exact lifts by the one-dimensional
space $\mathbf F_p\operatorname{id}_W$.
This gives precisely \eqref{eq:product-entire-fixed-space}.
The character projectors, their indicated combinations,
and the normalized top differential are defined over
a finite field extension when the original base is
finite; over any perfect base they descend after a
finite algebraic extension.  Over that extension the
whole geometric invariant space is already arithmetic.
\end{proof}

\begin{corollary}\label{cor:product-two-explicit-fixed-spaces}
For the explicit surface over $\mathbf F_3$ in
\eqref{eq:odd-product-explicit-elliptic}, the relative
rank is $9d^2-1$ and the full geometric and arithmetic
invariant covector spaces both have dimension two,
for every $d\geq1$.
For the surface over $\mathbf F_2$ in
Example~\ref{ex:ordinary-abelian-surface-explicit}, the
relative rank is also $9d^2-1$, and both invariant
covector spaces have dimension three for every even
$d\geq2$.
The exact full images inside the pointwise stabilizers
of these quotients are not determined by this theorem.
\end{corollary}
